\section{Institutions, by regime}
\label{sec:institutions}

This section draws out what follows from the measurements for institutions and for the people the
exposure ultimately reaches. It separates what the evidence supports from what it does not,
because recommendations are the easiest part of a paper of this kind to overreach.

\subsection{The organizing conclusion}

\begin{quote}
The federal government is fiscally exposed in both states, through different tax bases: wage taxes
if AI succeeds, capital gains and corporate taxes if it fails. Instruments that protect one base
do not protect the other.
\end{quote}

That is why the capital tax instrument is the only one that appears in both directions, and why
the contested base in Section~\ref{sec:fiscal} is the hinge rather than a technical footnote. The
rate links the two states without trading one against the other: as Section~\ref{sec:bust} shows,
bust-state revenue at risk is linear in the rate, so raising it scales both sides of the position
together, and the level cannot be sized without a measure of US AI capital income.

\subsection{What is needed at each regime}

For each regime we state the minimum instrument set identified by the measured mechanisms, which is a narrower claim than a policy recommendation. Appendix~\ref{app:detail} sets the four regimes and the failure case out in full. The shape of it is this. In the least-extrapolated scenario the set is the capital tax instrument, the trust fund contribution base, the student loan trigger and better measurement, and it is the only regime where we can say what is unnecessary as confidently as what is not: the housing agencies absorb that level without a Treasury draw and the banking system barely moves. At larger displacement nearly everything enters, with ownership and direct-claim instruments arriving precisely because the tax instrument is at its limit. Sudden displacement reorders the same set, since only contract instruments act inside a quarter, so pre-authorization is the whole instrument. At near-total displacement the required rate lies far outside observed experience and the question stops being financial regulation. In the failure case only concentration limits on AI-linked lending and the capital tax rate remain.


\subsection{What ownership would cost}

Because ownership is the instrument that survives at the extreme, its price is worth stating.
Raising the state's share of an AI side worth about \result{AISideTotalBn} billion dollars from
\result{AIFederalShareRounded} to ten percent would cost roughly \result{StakeForTenPctBn} billion
dollars and close about \result{GapClosedAtTenPct} percent of the gap, a stake about
\result{StakeVersusAgencyCapital} times the entire capital of the housing agencies. That is why
most instruments in Appendix~\ref{app:instruments} are not ownership instruments, and why the
ownership row cannot be dismissed: at near-total displacement it is the only row left, and the
cost of a position rises with the price of the asset.

\subsection{What each type of institution should change}

What each type of institution should change, tied to its measured exposure and paired with what the evidence does not support, is set out in Appendix~\ref{app:detail}.


\subsection{What the evidence does not support}

Two claims are tempting and neither survives the measurement: wholesale change to household mortgage underwriting at moderate displacement, and occupation-based underwriting. Appendix~\ref{app:detail} sets out the three measured grounds for the first and, for the second, the legal position, which is that occupation is not a prohibited basis under the Equal Credit Opportunity Act but that occupation-based underwriting in housing credit exposes a lender to a discriminatory-effects claim under the Fair Housing Act. Combined with the finding that the exposure is spread almost evenly across borrowers, the risk-adjusted case for it is weak and the case for portfolio-level disclosure is strong.


\subsection{Emerging markets}

Three facts change the answer outside a reserve-currency economy and they compound: imported AI capital accrues its taxable surplus abroad, a sovereign that cannot borrow in a currency it issues cannot absorb the shock, and many such economies are already on a steep debt path. The measurement itself is also unavailable there. Appendix~\ref{app:detail} gives the figures and the reason.


\subsection{State-contingent triggers}

Instruments are more useful attached to a trigger carrying its current value. The capital tax trigger is already tripped: the assembled rate is \result{TauKBarkai} percent against a threshold of \result{RequiredEasy} percent, both sides observed today. Prime-age nonemployment stands at \result{PrimeAgeNonemp} against the observed maximum of \result{PrimeAgeMax}, where the reemployment fit leaves the data, so crossing it is the moment the rest of this architecture can no longer be sized. Appendix~\ref{app:detail} gives the rest, and Appendix~\ref{app:instruments} the full instrument table.


\subsection{Implications for individuals, organisations and society}
\label{sec:implications}

Four consequences are worth stating plainly, because the exposure is ultimately borne by people.

\emph{For the people displacement reaches.} The thinnest buffers are one quintile up from the
bottom, among households carrying the commitments of those above them without the income to
service them through a gap in pay, so relief designed around the bottom quintile misses them, as
does relief designed around occupation. If displacement arrives through hiring that does not
happen, the private loss shifts onto people at the start of their working lives and becomes
invisible to the separations data most early-warning systems watch.

\emph{For managers.} Wage risk is a concentration to be governed rather than a credit category to
be underwritten: it is uneven in a way the aggregate hides, and it cannot be managed by tightening
household underwriting, because most of what arrives comes through the second round. The
governance question is whether an institution can state, in a form a board can read, how much of
its income and collateral depends on wages continuing to be paid, and what it has agreed in
advance to do if they are not.

\emph{Strategically.} Exposure to a fall in wage income sits on the federal balance sheet, and the
state's claim on the income that would replace those wages is the tax of
Section~\ref{sec:fiscal}, which falls short of what replacement needs. Acquiring a direct claim
instead is expensive, and a government that notices late will find the price of the asset has
risen for the same reason its revenue fell. An exposure spread almost evenly across the system
cannot be diversified away inside it, which makes pre-agreed rules and disclosure the instruments
with the best ratio of value to cost.

\emph{For society.} The claims wages pay are spread far more evenly than the wages themselves, so
a proportional fall in pay is regressive before any response is chosen. The same ownership
concentration runs through both halves of the paper: the capital tax reaches little because most
corporate claims sit where the income tax does not, and a bust transmits weakly because the assets
that would fall are held by the households least likely to change their spending. A statistic that
locates wage risk should be used to size public instruments, not to price individuals out of
credit for what they do for a living.
